\documentclass[12pt]{article}
\usepackage[utf8]{inputenc}
\usepackage[T1]{fontenc}
\usepackage{hyperref}
\title{Wadi Wadi Language Resources for K-12 Curriculum Development}
\author{Compiled by Claude Sonnet 5.5 (claude-sonnet-5-5) for Andr\'as Kornai}
\date{}
\begin{document}
\maketitle

\section{Introduction}
Wadi Wadi (ISO 639-3 xwd; also spelled Wathiwathi, Wotti-wotti, Piangil) is a dormant Australian language. AIATSIS locates it on the Murray River in Victoria and New South Wales, between a point above the Murrumbidgee junction and Swan Hill, at Piangil, and north to about Moolpa. Sources differ on its classification. Museums Victoria places it with the Western Kulin languages, and Wikipedia says Pama--Nyungan. AIATSIS also notes a separate Piangil record that may be the Weki Weki language. Searches found no dictionary, word list, primer or school programme specific to Wadi Wadi. Very little is publicly available.

\section{Primary Educational Resources}
No textbooks, primers or school materials for Wadi Wadi turned up in the searches. No school is known to teach it.

\section{Linguistic and Reference Materials}
\begin{itemize}
\item AIATSIS Austlang entry D4, Wadi Wadi (location, alternative names, references): \url{https://collection.aiatsis.gov.au/austlang/language/d4}
\item Wikipedia, Wadi Wadi language (overview, classification): \url{https://en.wikipedia.org/wiki/Wadi_Wadi_language}
\item Endangered Languages Project record: \url{https://endangeredlanguages.com/node/97773}
\item Museums Victoria, Victorian Aboriginal languages (overview): \url{https://museumsvictoria.com.au/bunjilaka/about-us/first-peoples/victorian-aboriginal-languages/}
\end{itemize}
The search summary also cited two works that I did not see as direct links. One is Ian Clark's 1996 report on Aboriginal language areas in Victoria, with feasibility plans for language retrieval. The other is Blake, Hercus, Morey and Ryan (2011), Pacific Linguistics 628, on the Mathi group. Both should be checked in library catalogues.

\section{Religious and Community Resources}
The search summary mentions the Wadi Wadi Land and Water Indigenous Corporation as one of several Traditional Owner groups. No direct link was found. Any use of language material should be agreed with the Traditional Owners.

\section{Digital and Audio Resources}
No audio, video or software resources for Wadi Wadi were found. The Victorian Collections search returned some items: \url{https://victoriancollections.net.au/collections?q=wadiwadi}

\section{Curriculum Development Context}
Wadi Wadi has no living speakers, and the surviving documentation is thin. A curriculum would be heritage and revitalisation oriented. It would depend on community governance and on archival material (AIATSIS, the Victorian Aboriginal Corporation for Languages). Related Kulin or Murray River languages with fuller documentation could serve as comparison resources, but only with community agreement.
\end{document}